\section{PHÂN TÍCH YÊU CẦU}

\subsection{Bài toán}
\hspace{1cm}
Olist là một sàn thương mại điện tử kết nối nhiều người bán nhỏ với khách hàng trên toàn
Brazil; hàng được người bán tự đóng gói và bàn giao cho đơn vị vận chuyển. Với quy mô lãnh
thổ lớn, giao hàng đúng hạn là yếu tố quyết định trải nghiệm khách hàng: trong dữ liệu,
6,8\% đơn đã giao bị trễ so với ngày hẹn, và 60,9\% số đơn trễ nhận đánh giá $\leq$ 2 sao.\\
Sàn cần: (1) theo dõi doanh số theo thời gian, ngành hàng và khu vực; (2) theo dõi hiệu suất
giao hàng và mức độ hài lòng theo tháng, theo tuyến vận chuyển; (3) phát hiện sớm các đơn có
nguy cơ giao trễ để can thiệp (liên hệ người bán, hãng vận chuyển, báo trước cho khách); và
(4) phát hiện sớm các khách hàng có nguy cơ đánh giá thấp để chăm sóc trước khi họ viết đánh
giá. Đây là bài toán điển hình đòi hỏi một kho dữ liệu tích hợp làm nền tảng phân tích, thay
vì truy vấn trực tiếp trên hệ thống OLTP vốn không tối ưu cho các truy vấn tổng hợp phức tạp.

\subsection{Người dùng cuối và vai trò}
\begin{table}[H]
\centering
\caption{Người dùng cuối và nhu cầu sử dụng hệ thống}
\begin{tabular}{@{}p{5cm}p{12cm}@{}}
\toprule
\textbf{Vai trò} & \textbf{Nhu cầu} \\
\midrule
Quản lý (Manager) & Xem doanh thu, hiệu suất sản phẩm/ngành hàng, tỷ lệ giao trễ và tỷ lệ
đánh giá thấp theo tháng, các tuyến giao trễ nhiều nhất thông qua \textit{Manager Dashboard}. \\
Vận hành / Chăm sóc khách hàng (Ops) & Xem danh sách đơn có rủi ro giao trễ và đơn có rủi ro
bị đánh giá thấp, xem đề xuất xử lý cho đơn giao trễ, gắn khách hàng rủi ro vào chiến dịch chăm sóc thông qua
\textit{Ops Console}; quản lý chiến dịch và xem khách đã gắn tại màn hình \textit{Campaign}. \\
\bottomrule
\end{tabular}
\end{table}

\subsection{Câu hỏi nghiệp vụ mà kho dữ liệu phải trả lời}
\begin{enumerate}
    \item Doanh thu và số lượng đơn hàng biến động thế nào theo ngày / tuần / tháng / quý,
          theo ngành hàng và theo khu vực?
    \item Sản phẩm, ngành hàng nào đang bán chạy hoặc bán chậm?
    \item Tỷ lệ giao trễ, thời gian giao và điểm đánh giá biến động thế nào theo tháng và theo
          tuyến (bang người bán $\to$ bang khách)? Giao trễ ảnh hưởng tới điểm đánh giá ra sao?
    \item Đơn đang vận chuyển nào có nguy cơ giao trễ hạn, nên ưu tiên can thiệp đơn nào trước,
          và nếu kéo dài ngày hẹn giao thì rủi ro giảm bao nhiêu?
    \item Đơn vừa giao nào có nguy cơ nhận đánh giá thấp, khách hàng nào nên được đưa vào chiến
          dịch chăm sóc?
\end{enumerate}
Các câu hỏi 1--3 được trả lời trực tiếp bằng truy vấn OLAP trên kho dữ liệu (Mục~3.3),
trong khi câu hỏi 4--5 đòi hỏi các mô hình dự đoán -- chính là các thành phần DSS được trình
bày ở Mục~\ref{sec:dss}.

\subsection{Nguồn dữ liệu}
\begin{itemize}
    \item \textbf{Dữ liệu lõi}: bộ dữ liệu mở \textit{Brazilian E-Commerce Public Dataset by
          Olist} trên Kaggle -- khoảng 99 nghìn đơn hàng (09/2016--10/2018), 112 nghìn dòng sản
          phẩm, 96 nghìn khách hàng, 3 nghìn người bán, 99 nghìn đánh giá và 1 triệu điểm tọa độ
          theo mã bưu chính. Dữ liệu ghi đầy đủ các mốc thời gian của đơn (mua, duyệt, bàn giao
          vận chuyển, giao tới khách, ngày hẹn giao), là cơ sở cho hai mô hình DSS.
    \item \textbf{Chiến dịch chăm sóc}: Olist không có thực thể này. Chiến dịch không phải dữ
          liệu nguồn và không đi qua ETL: Ops tự thêm, sửa, xóa chiến dịch trên giao diện, dữ liệu
          được lưu trong kho để phục vụ DSS. Danh sách này không dùng cho phân tích và không làm
          đầu vào cho mô hình DSS.
    \item Olist không có thông tin kênh bán, tồn kho và tên sản phẩm, nên các phân tích liên
          quan được thực hiện ở mức ngành hàng.
\end{itemize}

\subsection{Yêu cầu phi chức năng}
\begin{itemize}
    \item Kho dữ liệu phải hỗ trợ truy vấn tổng hợp (rollup, drill-down) với thời gian phản
          hồi phù hợp cho dashboard tương tác (dùng \textit{materialized view} tiền tính toán).
    \item Lịch sử thay đổi thuộc tính khách hàng (ví dụ địa chỉ nhận hàng) phải được lưu vết
          (Slowly Changing Dimension -- SCD) để phục vụ phân tích xu hướng theo thời gian.
    \item Dữ liệu nạp vào kho phải được kiểm tra chất lượng (data quality) trước khi phục vụ
          phân tích, tránh sai lệch kết quả báo cáo và mô hình DSS.
    \item Mô hình DSS chỉ được dùng thông tin đã biết tại thời điểm chấm điểm (không rò rỉ dữ
          liệu tương lai) và phải được đánh giá trên dữ liệu ở giai đoạn sau giai đoạn huấn luyện.
\end{itemize}
